\pdfoutput=1
\documentclass[11pt]{article}
\usepackage[T1]{fontenc}
\usepackage{lmodern,microtype}
\setlength{\emergencystretch}{1em}
\usepackage[letterpaper,margin=1in]{geometry}
\usepackage{amsmath,amssymb,amsthm,mathtools}
\usepackage[colorlinks=true,linkcolor=blue,citecolor=blue,urlcolor=blue]{hyperref}
\hypersetup{pdftitle={Sharp finite-memory laws for quantum Einstein--Rosen waves},pdfauthor={}}
\newtheorem{theorem}{Theorem}
\newtheorem{lemma}[theorem]{Lemma}
\theoremstyle{remark}
\newtheorem{remark}[theorem]{Remark}
\newcommand{\tr}{\operatorname{Tr}}
\newcommand{\Ad}{\operatorname{Ad}}
\newcommand{\TV}{\operatorname{TV}}
\newcommand{\dd}{\mathrm{d}}
\newcommand{\Id}{\mathrm{id}}
\newcommand{\norm}[1]{\left\lVert #1\right\rVert}
\newcommand{\ket}[1]{\lvert #1\rangle}
\newcommand{\bra}[1]{\langle #1\rvert}
\newcommand{\proj}[1]{\lvert #1\rangle\!\langle #1\rvert}
\title{Sharp finite-memory laws for quantum Einstein--Rosen waves}
\author{}
\date{}
\begin{document}
\maketitle
\begin{abstract}
The exact Hamiltonian for polarized Einstein--Rosen waves couples the
retained radial Fock modes to a conserved continuous spectral variable
of the discarded modes. What finite memory reproduces their entire
adaptive process, rather than just a one-time channel? For physical
interrogation time at most \(T\), retained spectral gap \(g_L\), high
spectral interval of width \(W\), and small operational error
\(\varepsilon\), we show that the optimal persistent memory dimension
is \(\Theta(1+A g_L T W/\sqrt{\varepsilon})\), where \(A\) is the
Hamiltonian scale. One positive classical grid attains this rate,
including for correlated initial states. The converse applies to
arbitrary quantum memories when the simulator is fixed across
preparations: a coherent retained-mode probe converts compression to
unitary programming, and a recycling decoder distinguishes a Fourier
code of preparation labels. Separately, one fixed smooth pure high
state forces dimension linear in interrogation time at fixed accuracy.
That lower bound uses two one-particle interrogations separated by an
actual retained-state replacement, without a tester quantum ancilla or
feedback. The split concerns radial modes, not spatial regions.
\end{abstract}

\section{Exact process and resource model}\label{sec:model}

In the Ashtekar--Pierri Fock quantization of polarized cylindrical
vacuum gravity \cite{AP}, the auxiliary Hamiltonian is the
second quantization of multiplication by radial momentum \(k>0\).
Split the one-particle modes at a fixed \(\Lambda>0\); their Fock
spaces are \(\mathcal H_L,\mathcal H_H\), with corresponding auxiliary
Hamiltonians \(H_{0,L},H_{0,H}\). In units \(\hbar=1\) write the exact
physical Hamiltonian and the low unitary at high spectral value \(r\) as
\begin{equation}\label{eq:ham}
 H=A(1-q\otimes R),\qquad
 q=e^{-bH_{0,L}},\quad R=e^{-bH_{0,H}},\qquad
 u_r(t)=e^{iAtrq},\qquad A,b>0 .
\end{equation}
The common scalar phase \(e^{-iAt}\) is dropped. All times here refer
to the asymptotic physical time, not the auxiliary free evolution
parameter. The apparent interaction in \eqref{eq:ham} comes from the
nonlinear relation between the boundary Hamiltonian and the auxiliary
energy. Nonlinear time correlations for this quantization were
studied in \cite{BGV}. A conserved inaccessible observable can
generate classical multi-time memory \cite{GRS}; here its continuous
spectrum and physically selectable normal preparations are essential.

For a normal initial joint state \(\rho\), let \(P_B=\mathbf1_B(R)\).
Its positive trace-class operator measure on the retained system is
\begin{equation}\label{eq:measure}
 M_\rho(B)=
 \tr_H\!\left[(1\otimes P_B)\rho(1\otimes P_B)\right].
\end{equation}
A tester's inert reference can be included in \(L\) in this partial
trace. For every retained instrument history with final event \(F\),
let \(F_r\) be its backward effect at the start, computed with
\(u_r\) in the free intervals. Then
\begin{equation}\label{eq:law}
 \Pr_\rho(F)=\int\tr\!\left[F_r\,M_\rho(\dd r)\right].
\end{equation}
Indeed every retained operation acts trivially on \(H\), and the
total unitary is diagonal in the spectral representation of \(R\).
For a step approximation of \(R\) the effect at the input is an
orthogonal spectral sum of the retained effects; taking a uniform
mesh limit gives \eqref{eq:law}. The sum in \eqref{eq:measure} is
countably additive in trace norm (orthogonality removes cross terms
after tracing \(H\)). The derivation uses positive instruments and
does not assume a density of \(M_\rho\), a sharp \(r\)-eigenvector,
or an actual measurement of the high sector. In the archive,
Paper~68 \cite{P68} gives the spectral process, causal-break memory
and exact unbounded Hankel rank. Papers~18 and 53 \cite{P18,P53}
concern positive and coherent sewing, respectively, in other
settings. Our resource question is approximation of the whole
continuous spectral process.

\paragraph{Preparations and controls.}
Fix a closed interval \(I=[r_-,r_+]\) of width \(W>0\) in the
non-vacuum continuous spectrum \(0<r<e^{-b\Lambda}\).
Allow every normal joint preparation with high support in \(I\).
A strict low cap \(0\le H_{0,L}\le E_L\) is optional; if imposed it
must hold at preparation and be preserved by each instrument outcome.
Set
\[
 g_L=\operatorname{diam}\operatorname{spec}(q|_L)
     =
 \begin{cases}
  1-e^{-bE_L},&\text{strict low cap},\\
  1,&\text{no low cap}.
 \end{cases}
\]
A tester has finitely many interrogation steps, arbitrary normal
retained instruments, an independent finite quantum reference, and
a branchwise total physical waiting time at most \(T\). These are
ideal mode operations; we do not charge apparatus complexity,
spectral resolution, or the duration of controls. Quantum
controlled operations at fixed clock times are allowed, but the total
free physical clock is common. The free waiting-time sum after a
classical outcome may depend on that outcome.

\paragraph{Simulation and persistent memory.}
A \emph{universal} simulator uses one fixed causal engine for all
these preparations and testers, with one CP initialization on the
joint state and no later access to that state except what it has
retained. Error is the supremum of total variation of the tester
outcome distributions. Equivalently one can optimize one final
event; for a one-output channel with a fresh arbitrary reference
this convention is half the diamond norm. A reset replaces the
accessible retained state by a fixed state and disposes of its old
state. We count in \(d\) all registers of the simulator carrying
information through a reset, including persistent classical
records and correlations in joint registers even if a marginal is
fixed. A hybrid memory's dimension is the sum over its classical
sectors. Independent, fixed resources prepared after the reset may
be absorbed into a common continuation; there is no other hidden
communication or selection of the engine based on the preparation.
A tester can withhold its recorded first outcome, and can insert
fresh data only after the reset.

\begin{theorem}[Universal finite-budget minimax law]\label{thm:univ}
For \(T>0\), \(g_L>0\), and sufficiently small absolute
\(0<\varepsilon\), there are universal constants \(c,C>0\) such that
the minimal persistent dimension in this model satisfies
\begin{equation}\label{eq:minimax}
 c\left(1+\frac{A g_LTW}{\sqrt{\varepsilon}}\right)
 \le d_{\mathrm{opt}}
 \le C\left(1+\frac{A g_LTW}{\sqrt{\varepsilon}}\right).
\end{equation}
The upper bound is attained in order by a fixed classical register
with a CP initialization. The lower bound permits arbitrary
mixed programs on a quantum memory and arbitrary fixed quantum
processing. An additional high-vacuum sector requires at most
one additional exact upper-bound label. At \(T=0\) or \(g_L=0\) no
hidden temporal memory is needed.
\end{theorem}

For a strict high total auxiliary-energy support cap \(E_H>\Lambda\),
the continuous width is \(e^{-b\Lambda}-e^{-bE_H}\). The same scaling
uses the full cap interval in the upper bound and any fixed interior
fraction in the lower bound. When the vacuum is admitted it is
separately represented. A known exact point law, if present as a
genuine spectral sector, also needs no varying program. The precision
exponent in \eqref{eq:minimax} does not concern a processor designed
for just one high preparation. Section~\ref{sec:fixed} gives a
separate fixed-preparation lower bound.

\section{Positive simulator for all adaptive histories}\label{sec:upper}

Choose a uniform grid \(r_j=r_-+jh\) in \(I\), with adjacent affine
weights \(w_j(r)\). For \(r\) between two grid points only their two
weights are nonzero, and
\[
 w_j(r)\ge0,\qquad
 \sum_j w_j(r)=1,\qquad
 \sum_j r_jw_j(r)=r.
\]
Initialize the retained space and a classical register by the
trace-preserving CP map
\begin{equation}\label{eq:init}
 \rho\longmapsto
 \sum_j\tr_H\!\left[(1\otimes \sqrt{w_j(R)})\rho
                      (1\otimes \sqrt{w_j(R)})\right]\otimes\proj j
   =\sum_j\left(\int w_j\,\dd M_\rho\right)\otimes\proj j.
\end{equation}
The retained marginal is unchanged. Apply \(u_{r_j}\) in the
\(j\)-sector during each dwell. There is one random label for the
entire run: choosing new labels at each cut need not preserve its
correlations. If included, the high vacuum is split off before
interpolation and represented at \(r=1\) exactly.

\begin{lemma}\label{lem:derivs}
For any tester in the model the backward effect satisfies
\[
 \norm{\partial_r F_r}\le A g_L T,\qquad
 \norm{\partial_r^2 F_r}\le (A g_L T)^2 .
\]
\end{lemma}
\begin{proof}
For a node with remaining time budget \(\tau\), dwell \(t\), and
outcome maps \(\mathcal I_o\), its effect is
\[
 \Ad(u_r(t))^*\left(\sum_o \mathcal I_o^*F_{o,r}\right).
\]
The flagged adjoint \((Y_o)\mapsto\sum_o\mathcal I_o^*Y_o\) is
unital and positive, so it contracts the supremum norm on
self-adjoint fields. Subtracting the spectral midpoint of \(q\)
changes no conjugation, and makes its operator norm at most \(g_L/2\).
One derivative of the conjugation costs at most \(A g_L t\) and two
at most \((A g_L t)^2\). Begin at an \(r\)-independent terminal event.
All child budgets are at most \(\tau-t\). The product rule then gives
\(A g_L[t+(\tau-t)]\) and
\((A g_L)^2[t^2+2t(\tau-t)+(\tau-t)^2]\), respectively. The same
binomial induction bounds the third derivative, giving uniform
Taylor remainders even for normal classical outcome kernels at
finite depth and measurable outcome-dependent dwells. The sum is
then the normal flagged integral and the adjoint contraction is
unchanged. An inert reference tensors the generator with identity.
\end{proof}

A norm-valued Taylor expansion at \(r\) annihilates the first-order
term in \(\sum_j w_j(r)F_{r_j}\). The second-order bound is
\((A g_LT)^2(r-r_j)(r_{j+1}-r)/2\) between successive grid points.
Pairing with the positive measure of total mass one in
\eqref{eq:law} gives
\begin{equation}\label{eq:errupper}
 \sup_{\mathrm{testers}}\TV(P_{\mathrm{true}},P_{\mathrm{grid}})
 \le \frac{(A g_L T h)^2}{8}.
\end{equation}
A uniform partition into
\(\max(1,\lceil A g_L T W/\sqrt{8\varepsilon}\rceil)\) intervals
therefore proves the upper bound. The second-order cancellation is
in the whole tested process, uniformly in its depth; first-order
cancellation by unitary mixing for single gates is not a new
principle \cite{Campbell}. Equation~\eqref{eq:errupper} handles an
operator-valued initial law without approximating or recovering a
density.

\section{Converse for a universal quantum simulator}\label{sec:lower}

\subsection{One physical dwell ladder realizes a gate family}

Pick a retained auxiliary energy \(E_*=E_L/2>0\) under the cap,
or finite \(E_*\) with \(1-e^{-bE_*}\ge 1/2\) without it. There is a
smooth normalized Fock packet \(v\) narrowly supported around that
total energy. For instance, symmetrize a product of sufficiently
many identical one-particle packets near \(E_*/n<\Lambda\).
Write \(a_*=1-e^{-bE_*}\) and \(\Omega=Aa_*\); in either case
\(\Omega\ge A g_L/2\). Let \(X\) swap \(v\) and the retained
vacuum and act trivially on their orthogonal complement.

After a fixed reset to the vacuum, introduce fresh data
\(S=\mathbb C^{K+1}\) with basis \(\ket j\).
Apply \(X\) coherently on every \(j>0\) branch at time zero, and
undo it on branch \(j>0\) at time \(t_j=jT/K\). At every time,
the clocks are common; only the retained control depends coherently
on \(j\). On ideal monochromatic packets, after discarding the
scratch retained mode, the gate is
\begin{equation}\label{eq:gate}
 V_r\ket j=e^{-i\alpha rj}\ket j,\qquad \alpha=\Omega T/K.
\end{equation}
This expression removes the common vacuum phase for the entire
time \(T\). If \(a(E)=1-e^{-bE}\) differs from \(a_*\) by at most
\(h_a\) across the actual support of \(v\), the operator norm error
of the data-output isometry is at most \(AT h_a\): on each
orthogonal control branch this bounds the vector error of the
uncomputed packet. The induced trace-norm channel error, with any
passive reference, is at most \(2AT h_a\). All steps preserve the
strict support cap when present, and no outcomes are postselected.

We use a finite calibration family of high one-particle packets.
At each designated \(r\), choose a smooth normalized packet around
energy \(-b^{-1}\log r>\Lambda\), with \(R\)-support of width as small
as needed. One-sided approximations at the endpoints of \(I\)
suffice. Its additional diamond error relative to \eqref{eq:gate}
is at most \(2\Omega T\sup |r'-r|\) over its actual support.
Choose the combined physical packet error at most \(2\varepsilon\);
the same low packet works for all calibrations and data.

Only a state \(\sigma_r\) of the counted \(d\)-dimensional memory
depends on the preparation after the reset. The low state and
fresh data are product with this memory, and the same fixed controls
are used in the continuation. Tracing all but the data output yields
a single CPTP processor \(\mathcal P\) on \(S\otimes\mathbb C^d\).
Since the tester permits arbitrary fresh data, an independent
reference and final joint effects, the uniform simulation guarantee
implies
\begin{equation}\label{eq:programm}
 \norm{\mathcal P(\,\cdot\,\otimes\sigma_r)-\Ad(V_r)}_\diamond
 \le\delta,\qquad \delta=4\varepsilon ,
\end{equation}
at all the finitely chosen calibration points. A simulator with a
fresh state-dependent continuation instead of one engine would not
imply \eqref{eq:programm}.

\subsection{Recycling one mixed program}\label{sec:recycle}

Exact no-programming already forbids storing arbitrarily many exact
unrelated gates in one fixed finite register \cite{NC}. Approximate
unitary programming and its connection to metrology are developed
in \cite{YRC}; qubit phase programming is an earlier example
\cite{VMC,ZB}. Here the needed statement concerns a finite interval
of phase rates and a gate ladder whose size grows with \(T\).
We give the recycling argument, including the mixed-program
and fixed-decoder details.

Let \(W_0:S\otimes\mathbb C^d\longrightarrow S\otimes E\) be a
fixed Stinespring isometry of \(\mathcal P\), with \(E\) enlarged
by a fixed unused sector if needed. Purify \(\sigma_r\) on a
register \(Q\) not used by the decoder, writing \(\ket{\psi_r}\).
Dimension-independent Stinespring continuity \cite{KSW}
in a common enlarged environment yields a unit vector
\(\ket{\phi_r}\in E\otimes Q\) such that
\begin{equation}\label{eq:ksw}
 \norm{(W_0\otimes I_Q)(\,\cdot\,\otimes\ket{\psi_r})
                 -V_r\,\cdot\,\otimes\ket{\phi_r}}_{\mathrm{op}}
 \le c_0\sqrt{\delta},
\end{equation}
with, for example, \(c_0=2\). The isometry bound implies distance in
full trace norm at most \(2c_0\sqrt\delta\) including a reference.

For a fixed state \(\tau\) on \(S\otimes\mathbb C^d\) define the
CPTP left inverse on the range of \(W_0\) by
\begin{equation}\label{eq:inverse}
 {\cal L}(Y)=W_0^*YW_0+
   \tr[(I-W_0W_0^*)Y]\,\tau.
\end{equation}
After the forward use of \(W_0\), retain the data output and apply
\({\cal L}\) to \(E\) and a \emph{fresh fixed dummy} data state
\(\eta\), not to the data output just obtained. On the ideal
environment \(\phi_r\), use \eqref{eq:ksw} evaluated on
\(V_r^*\eta V_r\) and contractivity of \({\cal L}\otimes\Id_Q\).
The recovered state is within \(2c_0\sqrt\delta\) in full trace norm
of \((V_r^*\eta V_r)\otimes\proj{\psi_r}\). Discard the dummy.
Including the original forward error, one use therefore applies
the gate and restores the program purification to within
\(4c_0\sqrt\delta\), also with previously retained data. The use
of \(V_r^*\) is exclusively in this error proof, not in the decoder.

Iterate the fixed decoder
\(m=\lfloor 1/(16c_0\sqrt\delta)\rfloor\) times on fresh data
registers, possibly entangled with each other. Telescoping makes
the total trace-norm error at most \(1/4\). For instance
\(0<\varepsilon\le2^{-16}\) gives
\(m\ge 1/(32c_0\sqrt\delta)\). As \(Q\) was never used, this is
a fixed CPTP decoder from the \emph{single mixed program} to the
data registers. The decoder is an information bound on the counted
memory, not a physical experiment that runs past time \(T\).

\subsection{Orthogonal packing resolves both time and error scales}

Take \(K=\max(1,\lceil\Omega T W/\pi\rceil)\); then
\(\alpha W\le\pi\). On \(S^{\otimes m}\) the attainable integer
weights of \eqref{eq:gate} are all \(s\) from \(0\) to \(mK\).
Pick one product basis vector \(\ket s\) for each weight and feed
their uniform superposition to the recycling decoder.
With \(L=mK+1\), its ideal output is
\[
 \ket{\eta_r}=L^{-1/2}\sum_{s=0}^{L-1}e^{-i\alpha rs}\ket s .
\]
Take \(r_j=r_-+2\pi j/(\alpha L)\) for
\(j=0,\ldots,N-1\), where
\begin{equation}\label{eq:N}
 N=1+\left\lfloor \frac{\alpha WL}{2\pi}\right\rfloor.
\end{equation}
The points are in \(I\), \(N\le L\), and the outputs are orthogonal
by the finite Fourier sum. These are exactly the finite
calibration points used in \eqref{eq:programm}. A fixed
measurement therefore distinguishes them with average probability
at least \(7/8\) after decoding. Pulling the measurement back gives
a POVM \(E_j\) on only one \(d\)-dimensional program. The bound
\[
 \frac1N\sum_j\tr(E_j\sigma_{r_j})
      \le \frac1N\sum_j\tr E_j=\frac dN
\]
uses \(0\le\sigma_{r_j}\le I\) and remains valid for arbitrary
mixed program states and fixed ancillary resources. Since
\(N\ge\max(1,m\Omega TW/(2\pi))\), this proves the lower bound in
Theorem~\ref{thm:univ}. In physical terms, the single dwell ladder
provides a frequency band of width \(\Omega T\); the extra factor
\(1/\sqrt\varepsilon\) enters through the information required to
program that band, not through additional physical evolution.

\section{Linear memory growth for one fixed smooth state}\label{sec:fixed}

The preceding converse varies high preparations. We next keep one
normal pure high state fixed for every \(T\). Neither a tester
quantum ancilla nor a state-dependent controller is needed.

Choose a fixed subinterval \(J\) with closure in \(I^\circ\) and
a period \(W_0< |J|\) that is a fixed fraction of \(W\), for example
\(W_0=W/2\) and \(|J|=3W/4\). Choose a real smooth bump
\(w\in C_c^\infty(J)\) positive on a subinterval of length greater
than \(W_0\). Set
\begin{equation}\label{eq:wrapped}
 B(r)=\sum_{n\in\mathbb Z}w(r+nW_0)^2,\qquad
 \xi(r)=\frac{w(r)}{\sqrt{W_0B(r)}},\qquad
 v_H(k)=\sqrt{be^{-bk}}\,\xi(e^{-bk}).
\end{equation}
The periodic denominator is smooth and strictly positive.
The high one-particle vector \(v_H\) is normalized and smooth with
compact support inside the allowed band. Its law
\(\mu(\dd r)=|\xi(r)|^2\dd r\) is exactly uniform after reduction
modulo \(W_0\), since \(\sum_n|\xi(r+nW_0)|^2=1/W_0\).

\begin{theorem}[One fixed preparation, two no-feedback windows]\label{thm:fixed}
For the single initial state \(\ket{0_L}\otimes v_H\) in
\eqref{eq:wrapped}, there is, for every \(N\ge2\) and \(\eta>0\),
a two-window retained experiment with no tester quantum ancilla,
total free time \(4\pi(N-1)/(AQW_0)\), and a genuine retained
reset between the windows, for which any simulator
accurate within \(\varepsilon\) in total variation satisfies
\begin{equation}\label{eq:fixed}
 \frac dN\ge \frac{2N^2+1}{3N^2}-3\eta^2-2\varepsilon .
\end{equation}
Here \(Q>0\) is the width of any fixed inner one-particle
\(q\)-band, independent of \(N\).
There is also a fixed finite-particle implementation with
\(Q=g_L/2\) and the same bound. Thus at fixed sufficiently small
accuracy this one preparation already requires
\(d=\Omega(1+A g_LTW)\) to simulate all allowed histories.
\end{theorem}

\begin{proof}
Let \(q_\ell=q_0+Q\ell/(N-1)\), \(0\le\ell<N\), in the chosen
inner band. For a normalized smooth bump \(\varphi_\Delta\) supported
in \((-\Delta,\Delta)\), use the normalized retained one-particle
vectors
\[
 v_\ell(k)=\sqrt{be^{-bk}}\,
                 \varphi_\Delta(e^{-bk}-q_\ell).
\]
Choose \(\Delta>0\) so the supports are disjoint and inside the
allowed band. Prepare \(v=N^{-1/2}\sum_\ell v_\ell\) and wait
\(T_N=2\pi(N-1)/(AQW_0)\). For fixed \(a\), define
\[
 r_j^{\,0}=a+j W_0/N,\qquad
 \psi_j=N^{-1/2}\sum_{\ell=0}^{N-1}
                    e^{iAT_Nr_j^{\,0}q_\ell}v_\ell ,
 \qquad j=0,\ldots,N-1 .
\]
The \(\psi_j\) form an orthonormal basis of the packet span
by the finite Fourier sum. Use the \emph{complete} POVM
\[
 E_j=\proj{\psi_j}+
        N^{-1}\left(I-\sum_{m=0}^{N-1}\proj{\psi_m}\right).
\]
The component outside the packet span is assigned uniformly, not
discarded. Because the packet profiles are identical translates in
\(q\), projection of \(u_r(T_N)v\) to that span has a common amplitude
\(\phi(AT_Nr)\), where
\(\phi(s)=\int|\varphi_\Delta(x)|^2 e^{isx}\dd x\).
The conditional probability is
\[
 p_j(r)=|\phi(AT_Nr)|^2\,p_j^0(r)
           +(1-|\phi(AT_Nr)|^2)/N,\qquad
 p_j^0(r)=\frac1{N^2}\left|\sum_{\ell=0}^{N-1}
   e^{2\pi i\ell((r-a)/W_0-j/N)}\right|^2 .
\]
For independent \(X,X'\) with density \(|\varphi_\Delta|^2\),
\(1-|\phi(s)|^2\le s^2\operatorname{Var}X\le s^2\Delta^2\),
by averaging \(1-\cos(s(X-X'))\). Taking
\(\Delta\le\eta/(AT_Nr_+)\) gives uniform variation at most
\(\eta^2\) from the ideal probability vector. No singular
spectral input is used.

Measure the first window and replace the low system by \(v\).
Concretely its terminal instrument can be
\(\rho_L\mapsto\tr(E_j\rho_L)\proj v\) on outcome \(j\).
Record that outcome, but give no feedback to the second window;
the second uses the same dwell and POVM. By \eqref{eq:law},
the exact joint probability is the mixture of \(p_x(r)p_y(r)\)
against the \emph{same} \(\mu\). The ideal first marginal is
uniform by \eqref{eq:wrapped}. Its ideal agreement probability is
\[
 a_N^{\,0}=\frac1{N^3}\sum_{|k|< N}(N-|k|)^2
          =\frac{2N^2+1}{3N^2}.
\]
To see the first equality, the Fourier coefficients of the modulus
squared of \(\sum_{\ell=0}^{N-1}e^{i\ell\theta}\) are \(N-|k|\).
Apply Parseval, the exact wrapped law and the sum over the \(N\)
equal-probability shifts. The true first marginal differs by at
most \(\eta^2\) in variation, and the true agreement by at most
\(2\eta^2\) by the product variation bound for each \(r\).

For the simulator, condition the normalized cut memory state
\(\sigma_j\) on its first label, with law \(\pi_j\).
Because the reset is fixed and the continuation receives no
first-label feedback outside memory, the second measurement is
one fixed POVM \(D_j\) on that memory. Its agreement satisfies
\[
 a_{\rm sim}=\sum_j\pi_j\tr(\sigma_jD_j)
       \le d/N+\TV(\pi,\mathrm{unif}),
\]
using \(\sigma_j\le I\) for the uniform weights and bounding the
change of weights by variation. Approximating the true joint law
bounds \(\TV(\pi,\mathrm{unif})\) by \(\varepsilon+\eta^2\)
and \(a_{\rm sim}\) below by \(a_N^0-2\eta^2-\varepsilon\).
This proves \eqref{eq:fixed}.

To obtain \(Q=g_L/2\) with a finite number of retained particles,
choose a compact \(q\)-band of that width in the full low spectral
interval. On its corresponding positive total-energy interval,
use \(k_i=E x_i\) with \(\sum_i x_i=1\), and a fixed
normalized symmetric simplex profile near \(x_i=1/n\). Take
\(n\) so large that every \(E x_i<\Lambda\) on its support. In
\(L^2(\prod_i\dd k_i)\) the factor \(E^{-(n-1)/2}\) cancels the
Jacobian of \((E,x)\), and \(\sqrt{be^{-bE}}\) changes the total
energy coordinate to \(q\). The resulting isometry from the
\(q\)-band intertwines \(q\) with multiplication and realizes the
same exact translated profiles with a fixed particle count.
Without a low cap an inner band of width \(1/2\) still corresponds
to finite total energy. Taking the largest admissible \(N\), and
using \(d\ge1\) at short times, proves the last claim.
\end{proof}

This converse is a two-window state-discrimination bound, not
an exact Hankel-rank test. In particular the bound is linear in
\(N\), not the square-root dimension bound from factoring a
classical correlation matrix through \(d^2\) operator entries.
At any fixed small \(\varepsilon>0\),
Theorem~\ref{thm:univ} supplies the matching time dependence
for this preparation, but not a matching precision law as
\(\varepsilon\downarrow0\) for a simulator tailored to it.

\section{Relation to prior results and scope}

Hashimoto, Yoshida and Murao \cite{HYM} already use unbiased
barycentric rounding for finite-dimensional smooth unitary
families (their Theorem~1 and Sec.~S3) and use mixed-program
recycling for the leading precision exponent (Sec.~S7). Both
mechanisms have antecedents in \cite{Campbell,YRC}. We do not
claim a new generic programming exponent. Here the upper bound
handles a correlated, operator-valued initial law and the full
infinite-dimensional adaptive process with constants uniform in
the number of cuts. Normal finite-energy Einstein--Rosen probes
realize the programming problem; exact Fourier discrimination
then gives uniform joint \(T,W,g_L,\varepsilon\) scaling in
dimension with no logarithmic loss. The fixed smooth preparation
converse is a different operational result, not inferred from
programming a varying family.

The persistent dimension quantifies simulation of a particular
physical family of mode histories. A reduced one-time channel can
agree with an effective noise model that uses fresh noise, while
such a model can miss the shared spectral record revealed by
the reset test. The underlying reason is conservation of the high
spectral variable, not exchange of high and low radial energy
during a free dwell. The bounds neither define spatial gravitational
subsystems nor include bandwidth restrictions or costs for physical
apparatus beyond the stated energy support and interrogation time.

\begin{thebibliography}{99}
\bibitem{AP}
A. Ashtekar and M. Pierri,
\emph{Probing Quantum Gravity Through Exactly Soluble Midi-Superspaces I},
J. Math. Phys. \textbf{37}, 6250--6270 (1996),
\href{https://arxiv.org/abs/gr-qc/9606085}{arXiv:gr-qc/9606085}.
\bibitem{BGV}
J. F. Barbero G., I. Garay and E. J. S. Villase\~nor,
\emph{Quantum Einstein--Rosen waves: Coherent states and \(n\)-point functions},
Class. Quantum Grav. \textbf{25}, 205013 (2008),
\href{https://arxiv.org/abs/0808.2561}{arXiv:0808.2561}.
\bibitem{GRS}
K. Goswami, A. K. Roy, V. Srivastava, B. Perez, C. Giarmatzi,
A. Gilchrist and F. Costa,
\emph{Hamiltonian characterisation of multi-time processes with classical memory},
New J. Phys. \textbf{27}, 114515 (2025),
\href{https://arxiv.org/abs/2412.01998}{arXiv:2412.01998}.
\bibitem{P68}
Quantum Gravity Research Archive,
\emph{Paper 68: What a Retained-Mode Observer Can Recover in Quantum Einstein--Rosen Dynamics},
\href{https://github.com/g3173028019-beep/ai-automated-quantum-gravity-research/tree/main/papers/68-einstein-rosen-spectral-memory}{source and manuscript}.
\bibitem{P18}
Quantum Gravity Research Archive,
\emph{Paper 18: Positive boundary sewing},
\href{https://github.com/g3173028019-beep/ai-automated-quantum-gravity-research/tree/main/papers/18-positive-boundary-sewing}{source and manuscript}.
\bibitem{P53}
Quantum Gravity Research Archive,
\emph{Paper 53: Coherent quantum shear sewing},
\href{https://github.com/g3173028019-beep/ai-automated-quantum-gravity-research/tree/main/papers/53-coherent-quantum-shear-sewing}{source and manuscript}.
\bibitem{Campbell}
E. Campbell,
\emph{Shorter gate sequences for quantum computing by mixing unitaries},
Phys. Rev. A \textbf{95}, 042306 (2017),
\href{https://arxiv.org/abs/1612.02689}{arXiv:1612.02689}.
\bibitem{NC}
M. A. Nielsen and I. L. Chuang,
\emph{Programmable Quantum Gate Arrays},
Phys. Rev. Lett. \textbf{79}, 321--324 (1997),
\href{https://arxiv.org/abs/quant-ph/9703032}{arXiv:quant-ph/9703032}.
\bibitem{KSW}
D. Kretschmann, D. Schlingemann and R. F. Werner,
\emph{The Information-Disturbance Tradeoff and the Continuity of
Stinespring's Representation},
IEEE Trans. Inf. Theory \textbf{54}, 1708--1717 (2008),
\href{https://arxiv.org/abs/quant-ph/0605009}{arXiv:quant-ph/0605009}.
\bibitem{YRC}
Y. Yang, R. Renner and G. Chiribella,
\emph{Optimal Universal Programming of Unitary Gates},
Phys. Rev. Lett. \textbf{125}, 210501 (2020),
\href{https://arxiv.org/abs/2007.10363}{arXiv:2007.10363}.
\bibitem{VMC}
G. Vidal, L. Masanes and J. I. Cirac,
\emph{Storing Quantum Dynamics in Quantum States: A Stochastic
Programmable Gate},
Phys. Rev. Lett. \textbf{88}, 047905 (2002),
\href{https://arxiv.org/abs/quant-ph/0102037}{arXiv:quant-ph/0102037}.
\bibitem{ZB}
M. Ziman and V. Bu\v{z}ek,
\emph{Universality and optimality of programmable quantum processors},
Acta Phys. Hung. B \textbf{26}, 277--291 (2006),
\href{https://arxiv.org/abs/quant-ph/0612218}{arXiv:quant-ph/0612218}.
\bibitem{HYM}
Y. Hashimoto, S. Yoshida and M. Murao,
\emph{Barycentric programming: Classical realization of programs for
universal quantum processors},
\href{https://arxiv.org/abs/2611.02441}{arXiv:2611.02441} (2026).
\end{thebibliography}
\end{document}
